\begin{table}[p]
\centering
\scriptsize
\caption[Identity-based bidder-exclusion and bidder-preference rules: representative global catalogue.]{\textit{Identity-based bidder-exclusion and bidder-preference rules: representative global catalogue.} Rules are organized by the mechanism through which they restrict the bidder pool. Full-exclusion rules (Block~A) generate the clean decomposition of this paper. Price-preference rules (Block~B) leave both bidder types active in the auction and therefore entangle the two margins. Target/liberalization rules (Block~C) operate through volumes or origins rather than per-auction bidding. SB = small business; SDVOSB = Service-Disabled Veteran-Owned SB; WOSB = Women-Owned SB.}
\label{tab:v3_external_policies}
\begin{threeparttable}
\begin{tabular}{p{4.6cm}p{2.6cm}p{3.2cm}p{3.6cm}}
\toprule
\textbf{Policy / authority} & \textbf{Jurisdiction (since)} & \textbf{Mechanism} & \textbf{Identity criterion} \\
\midrule
\multicolumn{4}{l}{\textbf{A. Full exclusion (this paper's decomposition applies cleanly)}} \\
SBA small-business set-aside (FAR 19.5) & US federal (1953) & Full set-aside: non-SB excluded from the contract & SB by NAICS size standard \\
8(a) Business Development program (13 CFR 124) & US federal (1978) & Full set-aside in certified contracts & Socially / economically disadvantaged SB \\
WOSB / EDWOSB (13 CFR 127) & US federal (2011) & Industry-gated full set-aside & Women-owned SB in under-represented NAICS industries \\
SDVOSB ``Rule of Two'' (38 USC 8127) & US federal (VA, 2006) & Full set-aside when $\geq 2$ SDVOSBs compete & Service-disabled veteran-owned SB \\
Berry Amendment (10 USC 4862) & US DoD (1941) & Full exclusion of non-US supply & US textiles, food, specialty metals \\
Brazilian Lei Complementar 123/2006 & Brazil federal (2006) & Full set-aside below R\$80k + subitem exclusivity & ME/EPP revenue threshold \\
Brazilian Law 14{,}133/2021 & Brazil federal (2021) & Full set-aside + regional origin preference & ME/EPP + regional origin \\
\textbf{S\~ao Paulo state SME-only rule} (Dec.\ 47.945/2003, 48.042/2003) & Brazil state (2003) & Full set-aside below R\$80k; extended to Group 65 in March 2018 & ME/EPP revenue threshold \\
\midrule
\multicolumn{4}{l}{\textbf{B. Price/score preference (partial: entangles intensive and extensive margins)}} \\
HUBZone program (13 CFR 126) & US federal (1997) & Set-aside \emph{plus} 10\% price preference against non-HUBZone bids & SB in historically underutilized business zone \\
Buy American Act (41 USC 8301) & US federal (1933) & 6\% / 12\% price preference for domestic end products & US manufactured end product \\
Italian SME preferences (D.Lgs.\ 50/2016) & Italy (2016) & Partial preferences and lot division favoring SMEs & SME per EC Recomm. 2003/361 \\
\midrule
\multicolumn{4}{l}{\textbf{C. Target / liberalization (operates through aggregate volumes or origins)}} \\
EU Directive 2014/24/EU (SME provisions) & European Union (2014) & Mandatory lot division and relaxed turnover thresholds & SME per EC Recomm. 2003/361 \\
WTO Government Procurement Agreement (GPA) & International (1994 revised) & Conditional national-treatment among signatories & Country origin (negative) \\
Trade Agreements Act (19 USC 2511) & US federal (1979) & Waives Buy American Act for GPA signatories & Country origin (negative) \\
\bottomrule
\end{tabular}
\begin{tablenotes}
\footnotesize
\item \textit{Notes:} Block A gathers rules that \emph{fully exclude} one bidder type from the auction. Under this regime the sheltered-bidding decomposition of Section~\ref{sec:model} applies as in the Brazilian context studied here: the intensive margin is recovered from drop-out bids within the restricted pool, and the entry margin is recovered from the observed post-policy participation rates. Block B gathers rules under which both types remain active but one receives a bidding advantage; the intensive and extensive margins cannot be separated without additional parametric structure, consistent with the partial-preference identification used in \citet{marion2007}, \citet{krasnokutskaya2011}, and \citet{atheyseira2011}. Block C gathers rules that operate through aggregate volume quotas, mandatory lot division, or country-origin liberalization rather than per-auction bidder eligibility. Global coverage is representative rather than exhaustive; the EU alone catalogues more than one hundred jurisdiction-level preference regimes \citep{bosio2022}.
\end{tablenotes}
\end{threeparttable}
\end{table}
